\begin{table*}[tbp]
\caption{Contrastes de TEOM, sistemas modificados e comparadores gravimétricos}\label{tab:teom}
\begingroup\singlespacing\fontsize{10}{12}\selectfont
\begin{tabularx}{\textwidth}{@{}P{.24\textwidth}P{.30\textwidth}X@{}}
\toprule\textbf{Relato e fração} & \textbf{Configuração, contexto e denominador} & \textbf{Resultado publicado e definição}\\\midrule
\textcite{rsl201}; \MPdoiscinco{} & Erfurt; TEOM a 50~\textcelsius{} versus Harvard Impactor. 32 dias sem aquecimento do impactador e 14 com aquecimento. & Razão entre médias TEOM/impactador: 0,74 e 1,06, respectivamente. As etapas ocorreram em períodos sucessivos.\\\addlinespace
\textcite{rsl144}; \MPdoiscinco{} & Nova York; TEOM convencional versus sistema SES. Addison: 843/715 pares; Queens: 383 conjuntos. & Diferença média automático--FRM em Addison: $-2{,}23/-1{,}22$~\ugm{}; em Queens: $-2{,}37/-0{,}93$~\ugm{}, para convencional/SES.\\\addlinespace
\textcite{rsl105}; \MPdoiscinco{} & FDMS versus FRM; Queens, 283 pares; Addison, 349 pares em 22 meses. & Queens: $\mathrm{FDMS}=1{,}25\,\mathrm{FRM}-0{,}63$, $R^2=0{,}95$. Addison: $\mathrm{FDMS}=1{,}088\,\mathrm{FRM}-0{,}004$, $R^2=0{,}945$. Interceptos em \ugm{}.\\\addlinespace
\textcite{rsl101}; \MPdoiscinco{} & TEOM diferencial versus FRM; Rubidoux. Resultados de 24~h; NC para a estimativa de médias. & Médias: 36,4 e 26,8~\ugm{}; diferença: 9,6~\ugm{}. A investigação química relaciona parte da diferença à perda de nitrato no FRM.\\\addlinespace
\textcite{rsl129}; \MPdez{} & TEOM versus referência gravimétrica SSI, Bakersfield. Até 61 dias de referência; número válido varia por instrumento. & Inclinação de regressão 0,37 e correlação $r=0{,}95$. A inclinação não foi tratada como estimativa de viés médio percentual.\\\addlinespace
\textcite{rsl106}; \MPdez{} & TEOM versus alto volume em dois confinamentos bovinos do Kansas; 208 e 98 pares. & Inclinações TEOM/alto volume: 1,07 ($R^2=0{,}95$) e 1,04 ($R^2=0{,}85$). Resultado condicionado aos sítios e ao aerossol avaliados.\\\bottomrule
\end{tabularx}
\Fonte{Elaboração própria a partir das estimativas e dos denominadores nos relatos citados.}
\NotaQuadro{TEOM: \textit{Tapered Element Oscillating Microbalance}; SES: \textit{Sample Equilibration System}; FDMS: \textit{Filter Dynamics Measurement System}; FRM: \textit{Federal Reference Method}; SSI: entrada seletiva por tamanho utilizada na referência do estudo. NC: denominador da estimativa não consolidado na extração focal, sem inferir ausência no original. Os contrastes não constituem observações independentes em uma escala comum.}
\endgroup
\end{table*}
